\documentclass[12pt]{article}

\usepackage[spanish]{babel}
\usepackage[utf8]{inputenc}
\usepackage{amsmath}
\usepackage{graphicx}
\usepackage{geometry}
\usepackage{float}
\usepackage{caption}
\usepackage{hyperref}

\geometry{a4paper, margin=2.5cm}

\title{Actividad de Aprendizaje 1.3 \\ 
Simulación de un Controlador Proporcional-Derivativo (PD) \\ 
Matemáticas para Ciencia de Datos 8SA 2026A}

\author{Emilio Izquierdo Montero}
\date{24 de febrero de 2026}

\begin{document}

\maketitle

\section{Introducción}

El presente trabajo consiste en la simulación de un controlador Proporcional-Derivativo (PD) aplicado al control de la velocidad de un vehículo. 
El objetivo es aplicar el concepto de derivada como tasa de cambio en un sistema dinámico modelado en tiempo discreto.

\section{Modelo Matemático del Sistema}

Se modela la dinámica del vehículo mediante la ecuación diferencial:

\begin{equation}
m \frac{dv}{dt} = u - b v
\end{equation}

Donde:

\begin{itemize}
\item $v$ = velocidad del vehículo (magnitud controlada)
\item $u$ = fuerza aplicada (señal de control)
\item $m$ = masa del vehículo
\item $b$ = coeficiente de fricción
\end{itemize}

Para implementar el sistema en tiempo discreto se utiliza una aproximación por diferencias finitas:

\begin{equation}
v[k+1] = v[k] + \Delta t \cdot \frac{1}{m} (u[k] - b v[k])
\end{equation}

\section{Controlador PD en Tiempo Discreto}

El error se define como:

\begin{equation}
e[k] = r[k] - v[k]
\end{equation}

La derivada numérica del error se aproxima como:

\begin{equation}
\frac{de}{dt} \approx \frac{e[k] - e[k-1]}{\Delta t}
\end{equation}

El controlador PD discreto se define como:

\begin{equation}
u[k] = K_p e[k] + K_d \frac{e[k] - e[k-1]}{\Delta t}
\end{equation}

Donde:

\begin{itemize}
\item $K_p$ = Ganancia proporcional
\item $K_d$ = Ganancia derivativa
\end{itemize}

\section{Parámetros Utilizados}

\begin{itemize}
\item Masa $m = 1000$ kg
\item Fricción $b = 50$
\item Tiempo de muestreo $\Delta t = 0.1$ s
\item Referencia $r = 20$ m/s
\item $K_p = 500$
\item $K_d = 100$
\end{itemize}

\section{Resultados}

\subsection{Velocidad del Sistema}

\begin{figure}[H]
\centering
\includegraphics[width=0.8\textwidth]{/home/emilio/proyectos/pd/velocidad_vs_tiempo.png}
\caption{Respuesta de la velocidad del sistema ante una referencia constante.}
\end{figure}

\subsection{Señal de Control}

\begin{figure}[H]
\centering
\includegraphics[width=0.8\textwidth]{/home/emilio/proyectos/pd/senal_de_control.png}
\caption{Señal de control generada por el controlador PD.}
\end{figure}

\section{Análisis del Comportamiento}

Se observa que:

\begin{itemize}
\item La acción proporcional ($K_p$) reduce el error rápidamente.
\item La acción derivativa ($K_d$) disminuye el sobreimpulso y mejora la estabilidad.
\item El sistema alcanza la referencia con un tiempo de establecimiento corto.
\end{itemize}

Si $K_p$ aumenta excesivamente, el sistema puede presentar oscilaciones.
Si $K_d$ es demasiado grande, puede amplificar el ruido y tambien puede generar inestabilidad.

\section{Conclusión}

El controlador PD permite regular eficazmente la velocidad del sistema modelado. 
La derivada actúa como un mecanismo de amortiguamiento que mejora la estabilidad del sistema dinámico en constante cambio.

Se comprobó la importancia del ajuste adecuado de las ganancias $K_p$ y $K_d$ para lograr un comportamiento estable y eficiente.

\end{document}
